\chapter{Conclusion}

Nebula implements a coherent simulation-in-the-loop framework for exploring a
PCIe Gen-2-oriented receiver equalizer. It connects parameterized SKY130 CTLE
benches, four-port channel processing, behavioral DFE analysis, staged failure
handling, caching, provenance, and multiple candidate-generation strategies.
The separation between optimization and evaluation is the project's most
important architectural contribution. The natural-language wrapper preserves
that separation: it translates explicit constraints and formats existing
pipeline output, while electrical measurements and verdicts remain outside the
LLM boundary.

Preserved real-SPICE artifacts demonstrate completed PPO training and feasible
receiver candidates. In the primary completed qualification study,
\code{pipeline\_ep0} passed 27/27 TT/SS/FF conditions at 512-bit candidate
fidelity on the public IEEE/IBM development channel. The result establishes the
value of explicit qualification but remains specific to the recorded channel,
model, reduced corner grid, fidelity, and finite stimulus.

The available comparison does not show PPO outperforming Random Search or CEM.
The fair single-seed trial produced one success for Random Search and none for
PPO or CEM. The defensible conclusion is therefore that PPO is integrated and
trainable, not that it is the best optimizer for this problem.

Nebula should be viewed as a well-instrumented Stage~1 research platform. A
physical DFE, complete PVT coverage, layout-aware devices, statistical
verification, and compliance-oriented validation are required before stronger
receiver-level claims can be made.
